% Bagean: metodhe
% Bab: induksi
% Pérangan: induksi-ing-N

\documentclass[../../../include/open-logic-section]{subfiles}

\begin{document}

\olfileid{mth}{ind}{inN}

\olsection{Induksi ing~$\Nat$}

Ing wujud kang paling prasaja, induksi iku tèknik
kanggo mbuktekaké asil tumrap kabèh bilangan asli.
Tèknik iki migunakaké kasunyatan yèn saka $0$,
kanthi nambah~$1$ bola-bali, kita bakal tekan
saben bilangan asli. Mula kanggo mbuktekaké
sawijining bab bener kanggo kabèh bilangan,
kita bisa (1)~nuduhaké yèn bab iku bener kanggo $0$,
lan (2)~nuduhaké yèn saben bener kanggo
bilangan~$n$, bab iku uga bener kanggo bilangan
sabanjuré~$n+1$. Yèn ukara ``bilangan~$n$
nduwèni sipat $P$'' dicekak dadi $P(n)$
(lan ``bilangan~$k$ nduwèni sipat $P$''
dadi $P(k)$, lan sapanunggalané), bukti
induksi yèn $P(n)$ kanggo kabèh $n \in \Nat$
kudu ngemot:
\begin{enumerate}
\item bukti $P(0)$, lan
\item bukti yèn kanggo sembarang~$k$,
  yèn $P(k)$ mula $P(k+1)$.
\end{enumerate}
Supaya cetha, anggepen kita wis nduwèni
(1) lan~(2). Saka (1), $P(0)$ bener.
Saka~(2), mliginé kita ngerti yèn
$P(0)$ mula $P(0+1)$, yaiku $P(1)$.
Iki katut saka pratelan umum ``kanggo
sembarang~$k$, yèn $P(k)$ mula
$P(k+1)$'' kanthi masang~$0$ ing panggonan~$k$.
Dadi, miturut modus ponens, kita olèh~$P(1)$.
Saka (2) maneh, kanthi masang~$1$ ing panggonan~$k$,
kita olèh: yèn $P(1)$ mula~$P(2)$.
Amarga~$P(1)$ wis kabuktèkaké, modus
ponens mènèhi~$P(2)$, lan saterusé.
Kanggo sembarang bilangan~$n$, sawisé
langkah iki diulang $n$~kaping, kita tekan~$P(n)$.
Dadi (1) lan~(2) bebarengan mbuktekaké~$P(n)$
kanggo saben $n
\in \Nat$.

Ayo deleng conto. Anggepen kita arep ngerti
pira cacah jumlah béda kang bisa metu saka
$n$ dadu. Sanajan katon anèh, miwitia
saka $0$ dadu. Tanpa dadu, mung ana siji
jumlah kang mungkin ``diuncalké'': ora ana
titik babar pisan, dadi jumlahé~$0$.
Mula cacah asil kang béda yaiku~$1$.
Yèn mung siji dadu, yaiku $n=1$,
ana enem nilai, saka $1$ nganti~$6$.
Kanthi rong dadu, kita bisa éntuk
saben jumlah saka $2$ nganti $12$,
yaiku $11$ kamungkinan. Kanthi telung
dadu, saben bilangan saka $3$ nganti
$18$ bisa metu, yaiku $16$ kamungkinan.
$1$, $6$, $11$, $16$: katon ana pola.
Mbokmenawa wangsulané $5n+1$?
Mesthi $5n+1$ iku cacah paling akèh
kamungkinan, amarga mung ana $5n+1$
bilangan saka $n$, nilai paling cilik
kang bisa metu saka $n$ dadu
(kabèh nuduhaké $1$), nganti $6n$,
nilai paling gedhé (kabèh nuduhaké $6$).

\begin{thm}
  Kanthi $n$ dadu, kabèh $5n+1$ nilai
  kang mungkin saka $n$ nganti $6n$
  bisa diéntukaké.
\end{thm}

\begin{proof}
Tetepna $P(n)$ minangka pratelan:
``Saben bilangan saka $n$ nganti~$6n$
bisa diéntukaké nganggo $n$~dadu.''
Kanggo migunakaké induksi, kita buktèkaké:
\begin{enumerate}
\item \emph{dhasar induksi} $P(1)$,
  yaiku kanthi siji dadu, saben bilangan
  saka $1$ nganti $6$ bisa metu.
\item \emph{langkah induksi}: kanggo
  saben $k$, yèn $P(k)$ mula~$P(k+1)$.
\end{enumerate}

Kasus tanpa dadu wis dibuktèkaké ing ndhuwur.
(1) Kabuktèkaké kanthi ndeleng dadu mawa
$6$ sisih. Dadu iku nduwèni enem sisih,
lan saben bilangan saka $1$ nganti~$6$
muncul ing salah sijiné. Mula kanthi siji
dadu saben bilangan saka $1$ nganti~$6$
bisa diéntukaké.

Kanggo mbuktekaké (2), kita nganggep
antesedèn kondisional, yaiku $P(k)$.
Anggepan iki diarani \emph{hipotesis induksi}.
Anggepan iki digunakaké kanggo mbuktekaké~$P(k+1)$.
Pérangan kang angel yaiku nemokaké cara
ndeleng nilai kang bisa metu saka $k+1$
dadu liwat nilai saka $k$~dadu lan nilai
dadu tambahan kang urutané $k+1$.
Nanging iki kudu ditindakaké supaya
hipotesis induksi bisa digunakaké.

Hipotesis induksi nyatakaké yèn kita
bisa éntuk saben bilangan saka $k$
nganti~$6k$ nganggo $k$~dadu.
Yèn kita éntuk~$1$ saka dadu urutan
$(k+1)$, jumlahé nambah $1$.
Dadi kita bisa éntuk saben nilai saka
$k+1$ nganti $6k+1$ kanthi nguncalaké
$k$~dadu banjur éntuk~$1$ saka
dadu urutan $(k+1)$. Nilai endi kang isih
durung? Nilai saka $6k+2$ nganti
$6k+6$. Iki bisa diéntukaké kanthi
$k$ dadu nuduhaké $6$s, banjur
salah siji bilangan saka $2$ nganti
$6$ metu saka dadu urutan $(k+1)$.
Mula kanthi $k+1$ dadu kita
bisa éntuk saben bilangan saka $k+1$
nganti $6(k+1)$, yaiku
wis mbuktekaké~$P(k+1)$ saka
anggepan~$P(k)$, hipotesis induksi.
\end{proof}

Kita asring migunakaké induksi kanggo
mbuktekaké bab sawijining rerangkening
obyèk, kaya bilangan utawa himpunan,
kang dhéwé didhéfinisikaké
``kanthi induktif'': obyèk urutan
$(n+1)$ didhéfinisikaké saka obyèk
urutan $n$. Upamané, jumlah~$s_n$
saka bilangan asli nganti~$n$ bisa
didhéfinisikaké kanthi
\begin{align*}
  s_0 & = 0\\
  s_{n+1} & = s_n + (n+1)
\end{align*}
Dhéfinisi iki mènèhi:
\begin{align*}
  s_0 & = 0,\\
  s_1 & = s_0 + 1 && = 1,\\
  s_2 & = s_1 + 2 && = 1 + 2 = 3\\
  s_3 & = s_2 + 3 && = 1 + 2 + 3 = 6, \text{ lsp.}
\end{align*}
Saiki kanthi induksi kita bisa
mbuktekaké yèn $s_n = n(n+1)/2$.

\begin{prop}
  $s_n = n(n+1)/2$.
\end{prop}

\begin{proof}
  Kita kudu mbuktekaké (1)
  $s_0 = 0\cdot(0 + 1)/2$ lan (2)
  yèn $s_k =
  k(k+1)/2$ mula
  $s_{k+1} = (k+1)(k+2)/2$.
  Pérangan (1) cetha. Kanggo
  mbuktekaké (2), kita nganggep
  hipotesis induksi:
  $s_k = k(k+1)/2$. Saka kéné
  kita kudu nuduhaké
  $s_{k+1} = (k+1)(k+2)/2$.

  Pira $s_{k+1}$? Miturut dhéfinisi,
  $s_{k+1} = s_k + (k+1)$.
  Miturut hipotesis induksi,
  $s_k = k(k+1)/2$. Iki bisa
  disulih menyang pepadhan sadurungé,
  banjur cukup ngétung pecahan:
  \begin{align*}
    s_{k+1} & = \frac{k(k+1)}{2} + (k+1) = {}\\
    & = \frac{k(k+1)}{2} + \frac{2(k+1)}{2} = {}\\
    & = \frac{k(k+1) + 2(k+1)}{2} = {}\\
    & = \frac{(k+2)(k+1)}{2}.
  \end{align*}
\end{proof}

Piwulang wigatiné: yèn arep mbuktekaké
pratelan babagan barisan $a_n$ kang
didhéfinisikaké kanthi induktif,
induksi dadi cara kang lumrah.
Senajan barisané ora didhéfinisikaké
mangkono, kaya kamungkinan asil
uncalan dadu, induksi tetep bisa
digunakaké yèn kasus~$k+1$
bisa digandhèngaké karo kasus~$k$.

\end{document}
